\section{Rank-one limits and a norm identity}
\label{sec:norm}

Let $D,E,B$ be the maps supplied by Proposition~\ref{prop:pencil}.
Thus $D$ and $E$ are positive maps with output sizes $a,c\geq1$ and
are \emph{unital}, meaning $D(I_4)=I_a$ and $E(I_4)=I_c$,
and $B(y)$ is an $a\times c$ matrix.
For unit vectors $u,v\in\R^4$, define the orthogonal projections
\begin{equation}
 P(v)=\operatorname{proj}_{\ker D(I-vv^\top)},
 \qquad
 R(u)=\operatorname{proj}_{\ker E(I-uu^\top)}.
 \label{eq:projections}
\end{equation}
These act on $\R^a$ and $\R^c$, respectively.
The next proposition recovers the scalar form~\eqref{eq:form} from the
Schur threshold~\eqref{eq:schur-equivalence}.

\begin{proposition}
\label{prop:norm}
Let $a,c\geq1$, let $D:\Sym^4\to\Sym^a$ and
$E:\Sym^4\to\Sym^c$ be positive unital linear maps, and let
$B:\R^3\to\R^{a\times c}$ be linear. Assume that
\eqref{eq:schur-equivalence} holds for every $X\succ0$,
$Z\in\Sym^4$ and $y\in\R^3$, with $\Phi_y$ defined
by~\eqref{eq:phi}. Define $P(v)$ and $R(u)$ by~\eqref{eq:projections}.
For every pair of unit vectors $u=(u_0,u'),v=(v_0,v')\in\R^4$
and every $y\in\R^3$,
\begin{equation}
 b(v_0u'-u_0v',y)
 =\norm{P(v)B(y)R(u)}_{\op}^{2}.
 \label{eq:norm-identity}
\end{equation}
Here the operator norm uses the Euclidean structures on the two block spaces.
\end{proposition}

\begin{proof}
Fix $u,v,y$. For $t\geq1$, set
\[
 X_t=vv^\top+t(I-vv^\top),
 \qquad
 Z_t=uu^\top+t(I-uu^\top).
\]
Both matrices are at least $I$. Positivity and unitality therefore give
$D(X_t)\succeq I$ and $E(Z_t)\succeq I$, so all inverse square roots below
exist. Put
\[
 A_t=Z_t^{-1/2}\Phi_y(X_t^{-1})Z_t^{-1/2},
 \qquad
 C_t=D(X_t)^{-1/2}B(y)E(Z_t)^{-1/2}.
\]
The matrix $A_t$ is positive semidefinite. For every real $r$, congruence
by $Z_t^{-1/2}$ gives
\[
 rZ_t\succeq\Phi_y(X_t^{-1})
 \quad\Longleftrightarrow\quad
 r\geq\lambda_{\max}(A_t).
\]
Likewise, linearity of $E$ and congruence by $E(Z_t)^{-1/2}$ give
\[
 \begin{split}
 E(rZ_t)\succeq B(y)^\top D(X_t)^{-1}B(y)
 &\quad\Longleftrightarrow\quad rI\succeq C_t^\top C_t\\
 &\quad\Longleftrightarrow\quad r\geq\norm{C_t}_{\op}^{2}.
 \end{split}
\]
Equation~\eqref{eq:schur-equivalence} identifies these two closed rays in
$r\in\R$, including their endpoints. Consequently,
\begin{equation}
 \lambda_{\max}(A_t)=\norm{C_t}_{\op}^{2}.
 \label{eq:threshold-norm}
\end{equation}
This applies to rectangular $C_t$ and includes a zero threshold.

The explicit formulas
\[
 X_t^{-1}=vv^\top+t^{-1}(I-vv^\top),
 \qquad
 Z_t^{-1/2}=uu^\top+t^{-1/2}(I-uu^\top)
\]
show that
\[
 A_t\longrightarrow
 uu^\top\Phi_y(vv^\top)uu^\top
 =\alpha uu^\top,
 \qquad
 \alpha=u^\top\Phi_y(vv^\top)u\geq0.
\]
The sign follows from positivity of $\Phi_y$. Thus the largest eigenvalue
of this limit is $\alpha$, also when $\alpha=0$. Continuity of the largest
eigenvalue gives $\lambda_{\max}(A_t)\to\alpha$.

For the other side, unitality gives
\[
 D(X_t)=I+(t-1)D(I-vv^\top).
\]
If $H\succeq0$ is a fixed finite-dimensional matrix, diagonalizing $H$
shows that
\begin{equation}
 [I+(t-1)H]^{-1/2}
 \longrightarrow\operatorname{proj}_{\ker H}
 \quad\text{in operator norm}.
 \label{eq:kernel-limit}
\end{equation}
Indeed, the factor on an eigenvector of eigenvalue $\mu\geq0$ is
$(1+(t-1)\mu)^{-1/2}$, which is one when $\mu=0$ and tends to zero
otherwise. Applying~\eqref{eq:kernel-limit} to $D(I-vv^\top)$ and
$E(I-uu^\top)$ yields
\[
 D(X_t)^{-1/2}\longrightarrow P(v),
 \qquad
 E(Z_t)^{-1/2}\longrightarrow R(u).
\]
Hence $C_t\to P(v)B(y)R(u)$ in operator norm. The dimensions $a,c$ are
fixed throughout; the limits include zero and identity kernel projections.
Taking limits in~\eqref{eq:threshold-norm} and using the rank-one
identity~\eqref{eq:rank-one} proves~\eqref{eq:norm-identity}.
\end{proof}

At $u=v$, the left side of~\eqref{eq:norm-identity} vanishes, so
$P(v)B(y)R(v)=0$ for every $y$.
We next study the identity as $u$ varies near $v$.
